\chapter{Information Bottleneck Training Loop}
\label{ch:ib_loop}

The Information Bottleneck (IB) principle provides a formal framework to understand and optimize deep neural networks by balancing compression and prediction. As derived in Chapter 10 (see Equation \ref{eq:ib_lagrangian}), the Information Bottleneck objective seeks a stochastic representation $Z$ of an input $X$ that maximally compresses $X$ while retaining maximum information about a target variable $Y$. The objective is formulated by minimizing the Information Bottleneck Lagrangian: $\mathcal{L}_{IB} = I(X; Z) - \beta I(Z; Y)$, where $\beta$ controls the tradeoff between compression of the input and preservation of relevant information for predicting the output. In practical training of deep neural networks, this translates to minimizing the prediction error (maximizing $I(Z; Y)$ via cross-entropy) while applying a penalty on the mutual information between the inputs and the intermediate representations (minimizing $I(X; Z)$). 

\begin{pycode}
import torch
import torch.nn as nn
import torch.optim as optim

def compute_mi_penalty(mu, logvar):
    """
    Computes a tractable upper bound on the mutual information I(X; Z).
    Following the variational approximation in Theorem \ref{thm:variational_ib},
    we use the Kullback-Leibler divergence between the posterior p(z|x) 
    and a prior q(z), often assumed to be a standard normal distribution.
    
    Args:
        mu (torch.Tensor): Mean of the latent representation.
        logvar (torch.Tensor): Log-variance of the latent representation.
        
    Returns:
        torch.Tensor: The KL divergence penalty representing I(X; Z).
    """
    # KL divergence between N(mu, diag(exp(logvar))) and N(0, I)
    # KL = -0.5 * sum(1 + log(sigma^2) - mu^2 - sigma^2)
    kl_div = -0.5 * torch.sum(1 + logvar - mu.pow(2) - logvar.exp(), dim=1)
    return kl_div.mean()
\end{pycode}

The code block above defines a function to compute a tractable upper bound on the mutual information term $I(X; Z)$. Direct computation of mutual information in high dimensions is typically intractable. Therefore, as established in Theorem \ref{thm:variational_ib}, we employ a variational approximation by introducing an auxiliary prior distribution $q(Z)$ and computing the Kullback-Leibler (KL) divergence between our encoder's representation $p(Z|X)$ and $q(Z)$. By assuming a Gaussian form for our encoder outputs (parameterized by a mean \texttt{mu} and a log-variance \texttt{logvar}) and a standard normal prior, the KL divergence resolves to a simple closed-form analytic expression. This penalty regularizes the network to discard irrelevant information from the inputs.

\begin{pycode}
def ib_training_loop(model, dataloader, beta=1e-3, epochs=10, lr=1e-3):
    """
    Minimal IB-regularized training loop.
    Minimizes: Cross-Entropy(Y, Y_hat) + beta * I(X; Z)
    """
    optimizer = optim.Adam(model.parameters(), lr=lr)
    criterion = nn.CrossEntropyLoss()
    
    for epoch in range(epochs):
        epoch_loss = 0.0
        epoch_ce = 0.0
        epoch_mi = 0.0
        
        for x, y in dataloader:
            optimizer.zero_grad()
            
            # Forward pass: model should return logits, mean, and log-variance
            logits, mu, logvar = model(x)
            
            # Calculate prediction loss (cross-entropy), corresponding to -I(Z; Y)
            ce_loss = criterion(logits, y)
            
            # Calculate compression penalty, corresponding to I(X; Z)
            mi_penalty = compute_mi_penalty(mu, logvar)
            
            # Full IB Lagrangian objective: L_IB = -I(Z;Y) + beta * I(X;Z)
            # which in practice is CrossEntropy + beta * MI_Penalty
            loss = ce_loss + beta * mi_penalty
            
            # Backward pass and optimization
            loss.backward()
            optimizer.step()
            
            epoch_loss += loss.item()
            epoch_ce += ce_loss.item()
            epoch_mi += mi_penalty.item()
            
        print(f"Epoch {epoch+1}/{epochs} | "
              f"Loss: {epoch_loss/len(dataloader):.4f} | "
              f"CE: {epoch_ce/len(dataloader):.4f} | "
              f"MI Penalty: {epoch_mi/len(dataloader):.4f}")
\end{pycode}

Here, we implement the complete training loop that minimizes the Information Bottleneck Lagrangian. The target neural network is expected to output not only the classification logits but also the parameters of the latent distribution (mean and log-variance) required to estimate the mutual information. For each batch, the algorithm computes the cross-entropy loss, which acts as a surrogate for maximizing $I(Z; Y)$, ensuring that the representation is highly predictive of the targets. Concurrently, it computes the mutual information penalty (the $I(X; Z)$ bound) and scales it by the Lagrange multiplier $\beta$. The sum of these two terms constitutes the final loss function that is backpropagated to update the model weights, seamlessly integrating information-theoretic regularization into standard deep learning pipelines.
